% ================================================================
%  Preambolo/tikz-pol.tex
%  Stili per le figure dell'appendice "Il poligono delle forze e il
%  poligono funicolare" (Volume III, A02).
%  Da richiamare con % ================================================================
%  Preambolo/tikz-pol.tex
%  Stili per le figure dell'appendice "Il poligono delle forze e il
%  poligono funicolare" (Volume III, A02).
%  Da richiamare con % ================================================================
%  Preambolo/tikz-pol.tex
%  Stili per le figure dell'appendice "Il poligono delle forze e il
%  poligono funicolare" (Volume III, A02).
%  Da richiamare con % ================================================================
%  Preambolo/tikz-pol.tex
%  Stili per le figure dell'appendice "Il poligono delle forze e il
%  poligono funicolare" (Volume III, A02).
%  Da richiamare con \input{Preambolo/tikz-pol} DOPO Preambolo/tikz.
%  Prefisso pol, sul modello degli stili fun di Preambolo/tikz-fun e
%  tcp di Preambolo/tikz-tcp; spessori e punte coincidono con quelli
%  di fun, perche' le figure dell'appendice sono richiamate dal
%  capitolo della fune e da quello della trave curva piana. Qui stanno
%  solo gli stili usati da piu' figure.
% ================================================================

\tikzset{
	pollato/.style={line width=1.1pt,line cap=round,line join=round},           % lati del poligono funicolare
	polforza/.style={-{Stealth[length=2.4mm,width=1.8mm]},line width=1.0pt},      % forze del sistema (nei due piani)
	polrisultante/.style={-{Stealth[length=2.8mm,width=2.1mm]},line width=1.5pt,black!55}, % risultante
	polestremo/.style={-{Stealth[length=2.4mm,width=1.8mm]},line width=1.0pt},    % forze di estremita' f_A, f_B
	polraggio/.style={line width=.5pt,line cap=round},                           % raggi polari
	polretta/.style={line width=.35pt,black!55},                                 % rette d'azione
	polcostruzione/.style={line width=.35pt,black!55,dash pattern=on 2pt off 1.5pt}, % prolungamenti, linee di costruzione
	polverso/.style={-{Stealth[length=1.5mm,width=1.1mm]},line width=.5pt},      % verso delle forze lungo un raggio o un lato
	polquota/.style={{Bar[width=2.4mm]Stealth[length=1.4mm,width=1.0mm]}-{Stealth[length=1.4mm,width=1.0mm]Bar[width=2.4mm]},line width=.4pt},
	polsuolo/.style={line width=.4pt,black!60},                                  % linee di riferimento
	% campo del piano di posizione: numero cerchiato (il polo e' il campo P)
	polcampo/.style={draw,circle,line width=.4pt,inner sep=0.5pt,minimum size=3.6mm,font=\scriptsize},
	% triangolo di equilibrio di un nodo, nel poligono delle forze
	poltriangolo/.style={fill=black!10},
}

% --- vertice del poligono delle forze, nodo del poligono funicolare, polo
\tikzset{
	pics/polpunto/.style={code={\fill (0,0) circle (1.2pt);}},
	pics/polpolo/.style={code={\fill (0,0) circle (1.6pt);}},
}

% --- lato di chiusura e poligono di prova
\tikzset{
	polchiusura/.style={line width=.7pt,dash pattern=on 4pt off 2pt},           % lato di chiusura
	polprova/.style={line width=.6pt,line cap=round,line join=round},            % poligono funicolare di prova
	polvincolo/.style={line width=.6pt,line join=round},                         % corpo dei vincoli
	polperno/.style={fill=white,draw=black,line width=.5pt},                     % perno delle cerniere
	polcarico/.style={-{Stealth[length=1.4mm,width=1.0mm]},line width=.4pt},     % forze distribuite
	polcurva/.style={line width=1.6pt,black!45,line cap=round},                  % curva funicolare
}

% --- vincoli a terra (dimensioni di funcerniera e funcarrello): perno nell'origine del pic
\tikzset{
	pics/polcerniera/.style={code={%
		\draw[polvincolo] (0,0) -- (-0.20,-0.34) -- (0.20,-0.34) -- cycle;
		\draw[polvincolo] (-0.36,-0.34) -- (0.36,-0.34);
		\foreach \k in {0,...,5}{
			\draw[line width=.35pt] ({-0.30+0.12*\k},-0.34) -- ++(-0.10,-0.12);}
		\draw[polperno] (0,0) circle (0.055);}},
	pics/polcarrello/.style={code={%
		\draw[polvincolo] (0,0) -- (-0.20,-0.30) -- (0.20,-0.30) -- cycle;
		\draw[line width=.5pt] (-0.10,-0.36) circle (0.06) (0.10,-0.36) circle (0.06);
		\draw[polvincolo] (-0.36,-0.42) -- (0.36,-0.42);
		\foreach \k in {0,...,5}{
			\draw[line width=.35pt] ({-0.30+0.12*\k},-0.42) -- ++(-0.10,-0.12);}
		\draw[polperno] (0,0) circle (0.055);}},
}
 DOPO Preambolo/tikz.
%  Prefisso pol, sul modello degli stili fun di Preambolo/tikz-fun e
%  tcp di Preambolo/tikz-tcp; spessori e punte coincidono con quelli
%  di fun, perche' le figure dell'appendice sono richiamate dal
%  capitolo della fune e da quello della trave curva piana. Qui stanno
%  solo gli stili usati da piu' figure.
% ================================================================

\tikzset{
	pollato/.style={line width=1.1pt,line cap=round,line join=round},           % lati del poligono funicolare
	polforza/.style={-{Stealth[length=2.4mm,width=1.8mm]},line width=1.0pt},      % forze del sistema (nei due piani)
	polrisultante/.style={-{Stealth[length=2.8mm,width=2.1mm]},line width=1.5pt,black!55}, % risultante
	polestremo/.style={-{Stealth[length=2.4mm,width=1.8mm]},line width=1.0pt},    % forze di estremita' f_A, f_B
	polraggio/.style={line width=.5pt,line cap=round},                           % raggi polari
	polretta/.style={line width=.35pt,black!55},                                 % rette d'azione
	polcostruzione/.style={line width=.35pt,black!55,dash pattern=on 2pt off 1.5pt}, % prolungamenti, linee di costruzione
	polverso/.style={-{Stealth[length=1.5mm,width=1.1mm]},line width=.5pt},      % verso delle forze lungo un raggio o un lato
	polquota/.style={{Bar[width=2.4mm]Stealth[length=1.4mm,width=1.0mm]}-{Stealth[length=1.4mm,width=1.0mm]Bar[width=2.4mm]},line width=.4pt},
	polsuolo/.style={line width=.4pt,black!60},                                  % linee di riferimento
	% campo del piano di posizione: numero cerchiato (il polo e' il campo P)
	polcampo/.style={draw,circle,line width=.4pt,inner sep=0.5pt,minimum size=3.6mm,font=\scriptsize},
	% triangolo di equilibrio di un nodo, nel poligono delle forze
	poltriangolo/.style={fill=black!10},
}

% --- vertice del poligono delle forze, nodo del poligono funicolare, polo
\tikzset{
	pics/polpunto/.style={code={\fill (0,0) circle (1.2pt);}},
	pics/polpolo/.style={code={\fill (0,0) circle (1.6pt);}},
}

% --- lato di chiusura e poligono di prova
\tikzset{
	polchiusura/.style={line width=.7pt,dash pattern=on 4pt off 2pt},           % lato di chiusura
	polprova/.style={line width=.6pt,line cap=round,line join=round},            % poligono funicolare di prova
	polvincolo/.style={line width=.6pt,line join=round},                         % corpo dei vincoli
	polperno/.style={fill=white,draw=black,line width=.5pt},                     % perno delle cerniere
	polcarico/.style={-{Stealth[length=1.4mm,width=1.0mm]},line width=.4pt},     % forze distribuite
	polcurva/.style={line width=1.6pt,black!45,line cap=round},                  % curva funicolare
}

% --- vincoli a terra (dimensioni di funcerniera e funcarrello): perno nell'origine del pic
\tikzset{
	pics/polcerniera/.style={code={%
		\draw[polvincolo] (0,0) -- (-0.20,-0.34) -- (0.20,-0.34) -- cycle;
		\draw[polvincolo] (-0.36,-0.34) -- (0.36,-0.34);
		\foreach \k in {0,...,5}{
			\draw[line width=.35pt] ({-0.30+0.12*\k},-0.34) -- ++(-0.10,-0.12);}
		\draw[polperno] (0,0) circle (0.055);}},
	pics/polcarrello/.style={code={%
		\draw[polvincolo] (0,0) -- (-0.20,-0.30) -- (0.20,-0.30) -- cycle;
		\draw[line width=.5pt] (-0.10,-0.36) circle (0.06) (0.10,-0.36) circle (0.06);
		\draw[polvincolo] (-0.36,-0.42) -- (0.36,-0.42);
		\foreach \k in {0,...,5}{
			\draw[line width=.35pt] ({-0.30+0.12*\k},-0.42) -- ++(-0.10,-0.12);}
		\draw[polperno] (0,0) circle (0.055);}},
}
 DOPO Preambolo/tikz.
%  Prefisso pol, sul modello degli stili fun di Preambolo/tikz-fun e
%  tcp di Preambolo/tikz-tcp; spessori e punte coincidono con quelli
%  di fun, perche' le figure dell'appendice sono richiamate dal
%  capitolo della fune e da quello della trave curva piana. Qui stanno
%  solo gli stili usati da piu' figure.
% ================================================================

\tikzset{
	pollato/.style={line width=1.1pt,line cap=round,line join=round},           % lati del poligono funicolare
	polforza/.style={-{Stealth[length=2.4mm,width=1.8mm]},line width=1.0pt},      % forze del sistema (nei due piani)
	polrisultante/.style={-{Stealth[length=2.8mm,width=2.1mm]},line width=1.5pt,black!55}, % risultante
	polestremo/.style={-{Stealth[length=2.4mm,width=1.8mm]},line width=1.0pt},    % forze di estremita' f_A, f_B
	polraggio/.style={line width=.5pt,line cap=round},                           % raggi polari
	polretta/.style={line width=.35pt,black!55},                                 % rette d'azione
	polcostruzione/.style={line width=.35pt,black!55,dash pattern=on 2pt off 1.5pt}, % prolungamenti, linee di costruzione
	polverso/.style={-{Stealth[length=1.5mm,width=1.1mm]},line width=.5pt},      % verso delle forze lungo un raggio o un lato
	polquota/.style={{Bar[width=2.4mm]Stealth[length=1.4mm,width=1.0mm]}-{Stealth[length=1.4mm,width=1.0mm]Bar[width=2.4mm]},line width=.4pt},
	polsuolo/.style={line width=.4pt,black!60},                                  % linee di riferimento
	% campo del piano di posizione: numero cerchiato (il polo e' il campo P)
	polcampo/.style={draw,circle,line width=.4pt,inner sep=0.5pt,minimum size=3.6mm,font=\scriptsize},
	% triangolo di equilibrio di un nodo, nel poligono delle forze
	poltriangolo/.style={fill=black!10},
}

% --- vertice del poligono delle forze, nodo del poligono funicolare, polo
\tikzset{
	pics/polpunto/.style={code={\fill (0,0) circle (1.2pt);}},
	pics/polpolo/.style={code={\fill (0,0) circle (1.6pt);}},
}

% --- lato di chiusura e poligono di prova
\tikzset{
	polchiusura/.style={line width=.7pt,dash pattern=on 4pt off 2pt},           % lato di chiusura
	polprova/.style={line width=.6pt,line cap=round,line join=round},            % poligono funicolare di prova
	polvincolo/.style={line width=.6pt,line join=round},                         % corpo dei vincoli
	polperno/.style={fill=white,draw=black,line width=.5pt},                     % perno delle cerniere
	polcarico/.style={-{Stealth[length=1.4mm,width=1.0mm]},line width=.4pt},     % forze distribuite
	polcurva/.style={line width=1.6pt,black!45,line cap=round},                  % curva funicolare
}

% --- vincoli a terra (dimensioni di funcerniera e funcarrello): perno nell'origine del pic
\tikzset{
	pics/polcerniera/.style={code={%
		\draw[polvincolo] (0,0) -- (-0.20,-0.34) -- (0.20,-0.34) -- cycle;
		\draw[polvincolo] (-0.36,-0.34) -- (0.36,-0.34);
		\foreach \k in {0,...,5}{
			\draw[line width=.35pt] ({-0.30+0.12*\k},-0.34) -- ++(-0.10,-0.12);}
		\draw[polperno] (0,0) circle (0.055);}},
	pics/polcarrello/.style={code={%
		\draw[polvincolo] (0,0) -- (-0.20,-0.30) -- (0.20,-0.30) -- cycle;
		\draw[line width=.5pt] (-0.10,-0.36) circle (0.06) (0.10,-0.36) circle (0.06);
		\draw[polvincolo] (-0.36,-0.42) -- (0.36,-0.42);
		\foreach \k in {0,...,5}{
			\draw[line width=.35pt] ({-0.30+0.12*\k},-0.42) -- ++(-0.10,-0.12);}
		\draw[polperno] (0,0) circle (0.055);}},
}
 DOPO Preambolo/tikz.
%  Prefisso pol, sul modello degli stili fun di Preambolo/tikz-fun e
%  tcp di Preambolo/tikz-tcp; spessori e punte coincidono con quelli
%  di fun, perche' le figure dell'appendice sono richiamate dal
%  capitolo della fune e da quello della trave curva piana. Qui stanno
%  solo gli stili usati da piu' figure.
% ================================================================

\tikzset{
	pollato/.style={line width=1.1pt,line cap=round,line join=round},           % lati del poligono funicolare
	polforza/.style={-{Stealth[length=2.4mm,width=1.8mm]},line width=1.0pt},      % forze del sistema (nei due piani)
	polrisultante/.style={-{Stealth[length=2.8mm,width=2.1mm]},line width=1.5pt,black!55}, % risultante
	polestremo/.style={-{Stealth[length=2.4mm,width=1.8mm]},line width=1.0pt},    % forze di estremita' f_A, f_B
	polraggio/.style={line width=.5pt,line cap=round},                           % raggi polari
	polretta/.style={line width=.35pt,black!55},                                 % rette d'azione
	polcostruzione/.style={line width=.35pt,black!55,dash pattern=on 2pt off 1.5pt}, % prolungamenti, linee di costruzione
	polverso/.style={-{Stealth[length=1.5mm,width=1.1mm]},line width=.5pt},      % verso delle forze lungo un raggio o un lato
	polquota/.style={{Bar[width=2.4mm]Stealth[length=1.4mm,width=1.0mm]}-{Stealth[length=1.4mm,width=1.0mm]Bar[width=2.4mm]},line width=.4pt},
	polsuolo/.style={line width=.4pt,black!60},                                  % linee di riferimento
	% campo del piano di posizione: numero cerchiato (il polo e' il campo P)
	polcampo/.style={draw,circle,line width=.4pt,inner sep=0.5pt,minimum size=3.6mm,font=\scriptsize},
	% triangolo di equilibrio di un nodo, nel poligono delle forze
	poltriangolo/.style={fill=black!10},
}

% --- vertice del poligono delle forze, nodo del poligono funicolare, polo
\tikzset{
	pics/polpunto/.style={code={\fill (0,0) circle (1.2pt);}},
	pics/polpolo/.style={code={\fill (0,0) circle (1.6pt);}},
}

% --- lato di chiusura e poligono di prova
\tikzset{
	polchiusura/.style={line width=.7pt,dash pattern=on 4pt off 2pt},           % lato di chiusura
	polprova/.style={line width=.6pt,line cap=round,line join=round},            % poligono funicolare di prova
	polvincolo/.style={line width=.6pt,line join=round},                         % corpo dei vincoli
	polperno/.style={fill=white,draw=black,line width=.5pt},                     % perno delle cerniere
	polcarico/.style={-{Stealth[length=1.4mm,width=1.0mm]},line width=.4pt},     % forze distribuite
	polcurva/.style={line width=1.6pt,black!45,line cap=round},                  % curva funicolare
}

% --- vincoli a terra (dimensioni di funcerniera e funcarrello): perno nell'origine del pic
\tikzset{
	pics/polcerniera/.style={code={%
		\draw[polvincolo] (0,0) -- (-0.20,-0.34) -- (0.20,-0.34) -- cycle;
		\draw[polvincolo] (-0.36,-0.34) -- (0.36,-0.34);
		\foreach \k in {0,...,5}{
			\draw[line width=.35pt] ({-0.30+0.12*\k},-0.34) -- ++(-0.10,-0.12);}
		\draw[polperno] (0,0) circle (0.055);}},
	pics/polcarrello/.style={code={%
		\draw[polvincolo] (0,0) -- (-0.20,-0.30) -- (0.20,-0.30) -- cycle;
		\draw[line width=.5pt] (-0.10,-0.36) circle (0.06) (0.10,-0.36) circle (0.06);
		\draw[polvincolo] (-0.36,-0.42) -- (0.36,-0.42);
		\foreach \k in {0,...,5}{
			\draw[line width=.35pt] ({-0.30+0.12*\k},-0.42) -- ++(-0.10,-0.12);}
		\draw[polperno] (0,0) circle (0.055);}},
}
